%==============================================================================
%  Shared body for the SIBD paper (class-agnostic).
%  Included by both sibd.tex (IEEEtran) and sibd_acm.tex (acmart) after \maketitle.
%  Title/author/abstract/keywords live in the class-specific driver files.
%==============================================================================

%==============================================================================
\section{Introduction}
\label{sec:intro}

Feature detection over a scale space---most famously the Scale-Invariant
Feature Transform (SIFT)~\cite{lowe2004}---remains a workhorse of computer
vision, underpinning registration, tracking, structure-from-motion, and
deformation measurement. Its computational profile, however, is a poor match
for the way it is usually implemented on GPUs. A canonical pipeline builds a
Gaussian pyramid, differences adjacent scales to form a Difference-of-Gaussians
(DoG) volume, searches that volume for local extrema, assigns orientations, and
emits descriptors. Each of these stages is typically realized as a separate
kernel that reads and writes large intermediate buffers to device memory. On
modern GPUs, where arithmetic throughput vastly outstrips off-chip memory
bandwidth, materializing the DoG volume and repeatedly streaming scale-space
images through DRAM is the dominant cost, not the convolutions or the
comparisons themselves.

This mismatch is especially acute for \emph{real-time video}, where a detector
must run every frame within a fixed latency budget and where per-frame memory
allocation and pyramid teardown add avoidable overhead. Existing GPU SIFT
implementations such as SiftGPU~\cite{siftgpu} and PopSift~\cite{popsift}
deliver strong throughput but retain the staged, memory-materializing
structure inherited from CPU designs, and they target the classic descriptor
rather than the specific needs of dense deformation measurement.

We present \textbf{SIBD}, a \emph{Scale-Invariant Blob Detector} that is
fully GPU-resident: the entire pipeline---from separable Gaussian smoothing to
a compacted list of characterized keypoints---executes without returning
control to the host and without ever writing a DoG volume to device memory.
SIBD makes three systems observations concrete. First, the DoG difference and
the $3\times3\times3$ scale-space extremum test can be \emph{fused in shared
memory}: a single kernel loads a small stack of adjacent Gaussian scales into
on-chip memory, forms the differences on the fly, performs the min/max
neighborhood test, and never spills the DoG. Second, separable convolution can
be tiled with overlapping \emph{aprons} so that halo regions are resolved
inside a thread block, and the column and row passes can be fused into a single
launch. Third, the sparse extrema can be packed into a dense keypoint list with
\emph{warp-aggregated} ballot-and-atomic compaction, avoiding a full-image
scan on the host.

Beyond systems structure, SIBD introduces a keypoint characterization tailored
to blobs. Rather than Lowe's histogram of gradient \emph{directions}, SIBD
accumulates an angular profile of the \emph{radial intensity derivative}
(a radial ``flux'') around each keypoint, binned by angular position. The
coefficient of variation of this profile is a direct, rotation-symmetric
measure of how blob-like a candidate is, and we use it to replace SIFT's
Hessian-based edge rejection with a single, cheaply computed symmetry filter.
This is well matched to our target application: tracking printed or painted
\emph{blob markers} on skin as an anatomical joint flexes, and recovering a
dense \emph{strain field} from the resulting position field.

\vspace{0.5em}
\noindent\textbf{Contributions.} We claim the following:
\begin{itemize}
  \item A \emph{fully GPU-resident} scale-space blob detector that emits a
        compacted keypoint list with \emph{zero} intermediate DoG volume in
        DRAM, via shared-memory fusion of DoG and extrema detection.
  \item An apron-tiled separable Gaussian pyramid with an optional
        single-launch fused column/row convolution, and warp-aggregated stream
        compaction of sparse keypoints.
  \item A novel \emph{radial-flux angular signature} and a
        \emph{coefficient-of-variation symmetry filter} for keypoint selection,
        replacing Hessian edge rejection.
  \item Real-time, multi-frame operation via \emph{persistent device-buffer
        reuse}, and an end-to-end application to strain-field capture from
        blob-marked skin over a flexing joint.
  \item An evaluation against SIFT, SiftGPU, PopSift, and OpenCV-CUDA on
        standard repeatability benchmarks and on strain ground truth.
\end{itemize}

\begin{fillbox}
Add a teaser figure (Fig.~1): input frame of blob-marked skin $\rightarrow$
SIBD keypoints $\rightarrow$ recovered strain field, with a latency/throughput
callout. Add one-sentence quantitative headline once results exist.
\end{fillbox}

%==============================================================================
\section{Related Work}
\label{sec:related}

\subsection{Scale-space feature detection}
SIFT~\cite{lowe2004} established the Gaussian/DoG scale space, extremum
localization, orientation assignment, and gradient-histogram descriptor that
most subsequent detectors refine. SURF~\cite{surf} approximates the Laplacian
with box filters and integral images for speed; KAZE/A-KAZE~\cite{kaze} move to
nonlinear diffusion scale spaces. Blob detectors based on the Laplacian- or
Hessian-of-Gaussian and their affine-covariant variants~\cite{mikolajczyk2005}
are the closest classical relatives of SIBD, since we detect and characterize
blobs rather than corners.

\subsection{GPU implementations of SIFT}
SiftGPU~\cite{siftgpu} pioneered a GPU pyramid and DoG search using graphics
pipelines; PopSift~\cite{popsift} and CudaSift~\cite{cudasift} provide modern
CUDA implementations with competitive throughput, and OpenCV's CUDA module
ships a production SIFT. These systems retain the staged structure of the CPU
algorithm: the Gaussian pyramid and DoG volume are materialized in device
memory and consumed by later kernels.
\begin{fillbox}
State precisely how each baseline handles the DoG volume and orientation stage,
and cite reported throughput numbers we will compare against. Note which ones
support persistent/streaming operation for video.
\end{fillbox}

\subsection{Kernel fusion and on-chip data reuse}
The bandwidth-bound nature of stencil and image pipelines motivates tiling with
halos/aprons and fusing producer--consumer stages to keep intermediates on
chip~\cite{halide,micikevicius_stencil}. Warp-aggregated atomics and ballot-based
stream compaction are established primitives for packing sparse
outputs~\cite{cuda_warpagg}.
\begin{fillbox}
Position SIBD's DoG+extrema fusion relative to generic stencil-fusion work:
our fusion is across the \emph{scale} dimension (a stack of Gaussian layers),
which is what removes the DoG volume entirely.
\end{fillbox}

\subsection{Marker-based deformation and strain measurement}
Digital image correlation (DIC) and marker tracking recover displacement and
strain fields from surface patterns~\cite{dic}. On-skin strain sensing for
biomechanics uses printed patterns or fiducials tracked across frames.
\begin{fillbox}
Contrast SIBD (sparse, scale-invariant blob markers tracked in real time on the
GPU) with dense DIC (subset correlation, typically offline). Argue why
scale-invariant blob markers are robust to the strong local scale changes that
occur when skin stretches over a flexing joint.
\end{fillbox}

%==============================================================================
\section{Background and Problem Setup}
\label{sec:background}
\subsection{Gaussian scale space and DoG extrema}
\begin{fillbox}
Define the Gaussian pyramid, octaves/layers, $\sigma_0$, $k$, and the DoG
approximation to the scale-normalized Laplacian. State the $3\times3\times3$
extremum criterion. Tie symbols to code: \texttt{sigma0}, \texttt{k},
\texttt{octaves}, \texttt{layers=scales}.
\end{fillbox}

\subsection{The strain-field capture problem}
\begin{fillbox}
Formalize: given blob markers on skin, detect and match blobs across frames to
get a position field $\mathbf{u}(x,y,t)$; derive the strain tensor from spatial
derivatives of the displacement field. Define accuracy metrics vs.\ ground
truth (mechanical rig / synthetic warp).
\end{fillbox}

%==============================================================================
\section{The SIBD Pipeline}
\label{sec:pipeline}
\begin{fillbox}
Overview figure: dataflow diagram showing which buffers live only in shared
memory (DoG, extrema) vs.\ device memory (Gaussian layers, keypoint list).
Emphasize the absence of a DoG volume in DRAM.
\end{fillbox}

\subsection{Apron-Tiled Separable Gaussian Pyramid}
\label{sec:gauss}
\begin{fillbox}
Describe \texttt{col\_kernel\_strips} and \texttt{row\_kernel}: strip/image
indexing, apron sizing per layer (\texttt{apron = ceil(3*sigma/2)*2}), shared
memory layout, and boundary handling. Present the optional fused
\texttt{col\_row\_kernel} (single launch) and when it wins. Include the block/
thread configuration table from the tuning comments (256 col threads,
$32\times12$ row threads, \texttt{maxregs}).
\end{fillbox}

\subsection{Fused DoG + Scale-Space Extrema (No DRAM DoG)}
\label{sec:fused}
\begin{fillbox}
Describe \texttt{blobs\_2}: load 5 adjacent Gaussian layers into 4 shared-memory
DoG planes, compute differences in place, run the flattened min/max 26-neighbor
test across layers, write only extrema. Quantify DRAM traffic saved vs.\
materialized DoG. Include the shared-memory budget calculation.
\end{fillbox}

\subsection{Warp-Aggregated Keypoint Compaction}
\label{sec:compact}
\begin{fillbox}
Describe \texttt{stream\_compact}: per-warp \texttt{vote\_ballot\_sync} mask,
\texttt{popc} counts, warp-leader atomic into block offset, block-leader atomic
into global counter, contrast threshold $\tau=0.01$. Emit \texttt{potential\_blob}
records (x, y, image, octave, layer).
\end{fillbox}

\subsection{Radial-Flux Angular Signature}
\label{sec:orient}

The kernel \texttt{find\_orientations} characterizes each candidate not by a
histogram of gradient \emph{directions}, as in Lowe's orientation
assignment~\cite{lowe2004}, but by an angular profile of the \emph{radial
intensity derivative} around the keypoint. One thread block processes one
keypoint; a $(2r{+}1)\times(2r{+}1)$ neighborhood (plus a one-pixel apron) at
the detected scale is staged in shared memory, where $r$ is the blob radius
determined at detection ($r=\lceil 1.5\,\sigma_0 k^{\ell}+1\rceil$ in code).

Let the keypoint center be $\mathbf{c}=(X,Y)$ and a neighborhood pixel be
$\mathbf{p}=(x,y)$, with displacement $\mathbf{d}=\mathbf{p}-\mathbf{c}$, radius
$\rho=\lVert\mathbf{d}\rVert$, and inward radial unit vector
$\hat{\mathbf{r}}=-\mathbf{d}/\rho$. The image gradient
$\nabla I(\mathbf{p})=(I_x,I_y)$ is estimated by central differences. Following
the Oral-Qual formulation, the neighborhood accumulates a projection of the
gradient onto the radial direction. We state the \emph{code-faithful} form,
which augments Eq.~(1) of the Oral-Qual writeup in three respects
(marked in \textcolor{teal}{teal}).

\vspace{0.3em}
\noindent\textbf{Per-pixel radial response.}
\begin{equation}
\label{eq:radial}
m(\mathbf{p}) \;=\; \nabla I(\mathbf{p})\cdot\hat{\mathbf{r}}
\;=\; -\,\frac{\nabla I(\mathbf{p})\cdot\mathbf{d}}{\rho},
\qquad \rho>0 .
\end{equation}
This is the (inward) \enrich{radial \emph{derivative}} of intensity---a
projection onto the \enrich{\emph{unit}} radial direction---rather than the raw
dot product with the full position vector $\mathbf{r}$ written in Oral-Qual
Eq.~(1). \enrich{At the singular center pixel ($\rho=0$) the kernel instead uses
$m=\tfrac12\lVert\nabla I\rVert$ and bins by the gradient angle}, a case the
original text did not treat.

\vspace{0.3em}
\noindent\textbf{Gaussian window.}
\begin{equation}
\label{eq:weight}
w(\mathbf{p}) \;=\; \frac{1}{2\pi r}\,
\exp\!\left(-\frac{\rho^{2}}{2r^{2}}\right).
\end{equation}
\enrich{This spatial weighting, present in the kernel but omitted from Oral-Qual
Eq.~(1), down-weights peripheral pixels and makes the signature robust to
neighborhood-boundary effects.}

\vspace{0.3em}
\noindent\textbf{Angular signature.} Each pixel contributes to the bin indexed
by its \emph{angular position} about the center,
\begin{equation}
\label{eq:bin}
\theta(\mathbf{p}) = \operatorname{atan2}(Y-y,\,X-x),\qquad
b(\mathbf{p}) = \Big\lfloor \tfrac{\theta(\mathbf{p})\bmod 2\pi}{2\pi/B}\Big\rfloor,
\end{equation}
producing the $B$-bin signature ($B=32$, the warp width)
\begin{equation}
\label{eq:signature}
h[b] \;=\!\! \sum_{\substack{\mathbf{p}\in N\\ b(\mathbf{p})=b}}\!\! w(\mathbf{p})\,m(\mathbf{p}).
\end{equation}
\enrich{We emphasize that the binning variable is the pixel's angular
\emph{position}, $\theta(\mathbf{p})$, not the gradient's direction; the
Oral-Qual prose was ambiguous on this point, and the two are distinct.}

\begin{enrichbox}
\textbf{Rotation invariance (added argument).} Rotate the image about
$\mathbf{c}$ by $\phi$. Positions and gradients co-rotate, so $\hat{\mathbf{r}}
\mapsto R_\phi\hat{\mathbf{r}}$ and $\nabla I\mapsto R_\phi\nabla I$; the
projection $m=\nabla I\cdot\hat{\mathbf{r}}$ is a dot product of two co-rotating
vectors and is therefore \emph{invariant}. The weight $w$ depends only on
$\rho$ and is invariant. The bin index shifts uniformly, $\theta\mapsto
\theta+\phi$, so the signature undergoes a cyclic shift, $h[b]\mapsto
h[b-\Delta]$. Consequently any \emph{shift-invariant} summary of $h$---in
particular its mean and standard deviation, hence the coefficient of variation
of \S\ref{sec:cov}---is exactly rotation invariant (up to bin quantization).
Note this yields a rotation-invariant \emph{selection criterion}, not a
canonical orientation as in Lowe's method.
\end{enrichbox}

\begin{fillbox}
Figure: the signature $h[b]$ for (a) an ideal blob (flat, low variation),
(b) an edge (bimodal, high variation), (c) a skin wrinkle. Add a note verifying
the central-difference sign convention against the shared-memory index stride
in \texttt{find\_orientations} (\texttt{dy = data[l-1]-data[l+1]}).
\end{fillbox}

\subsection{Coefficient-of-Variation Symmetry Filter}
\label{sec:cov}

A candidate is accepted only if its angular signature $h[\cdot]$ is nearly
uniform, which holds for radially symmetric blobs and fails for edges,
wrinkles, and other anisotropic structure. We quantify uniformity by the
coefficient of variation of the signature (Oral-Qual Eq.~(2), verified against
\texttt{filter\_blobs}):
\begin{equation}
\label{eq:cov}
c \;=\; \frac{\sigma(h)}{\lvert\mu(h)\rvert},\qquad
\mu(h)=\frac{1}{B}\sum_{b=1}^{B} h[b],\quad
\sigma(h)=\sqrt{\frac{1}{B}\sum_{b=1}^{B}\big(h[b]-\mu(h)\big)^2},
\end{equation}
and retain the keypoint iff $c<\theta_{\mathrm{cv}}$ (the code uses
$\theta_{\mathrm{cv}}\approx 0.6$). \enrich{We take the absolute value of the
mean and guard the near-zero case ($\lvert\mu\rvert<\epsilon\Rightarrow
c=\infty$), matching the kernel; the Oral-Qual Eq.~(2) wrote $\sigma/\mu$
without these safeguards.} Because $c$ is a shift-invariant functional of $h$,
the accept/reject decision inherits the rotation invariance argued in
\S\ref{sec:orient}.

This single scalar test replaces SIFT's two-part refinement---the low-contrast
threshold on the DoG response and the Hessian principal-curvature ratio
$\mathrm{tr}(H)^2/\det(H)$ used for edge rejection~\cite{lowe2004}. \enrich{In
contrast to the Hessian test, which requires forming and factoring a
$2\times2$ second-moment matrix per candidate, the CV test reuses the already
computed signature and costs one warp-level mean/variance reduction}, mapping
naturally to \texttt{shfl\_down\_sync} across the $B=32$ bins.

\begin{fillbox}
Threshold sweep $\theta_{\mathrm{cv}}$ vs.\ precision/recall on blob markers
(ablation, \S\ref{sec:results}). Compare rejection behavior against the Hessian
edge filter on the same candidates. Add the warp-reduction cost model
($O(\log B)$ shuffles) and cite Fig.~7 analogue.
\end{fillbox}

\subsection{Persistent Buffers for Real-Time Video}
\label{sec:persist}
\begin{fillbox}
Describe the \texttt{reset}/\texttt{gpu\_elements} reuse path in the newer
driver: allocate pyramid/keypoint buffers once, zero-and-reuse each frame,
amortize allocation. Report per-frame latency with vs.\ without reuse.
\end{fillbox}

%==============================================================================
\section{Implementation}
\label{sec:impl}
\begin{fillbox}
Julia + CUDA.jl. Kernel launch configs, \texttt{maxregs} choices, occupancy,
shared-memory sizing. Note portability considerations and any autotuning.
\end{fillbox}

%==============================================================================
\section{Experimental Setup}
\label{sec:setup}
\begin{fillbox}
GPU(s) and driver/toolkit versions. Baselines: reference SIFT (VLFeat/Lowe),
SiftGPU, PopSift, CudaSift, OpenCV-CUDA. Datasets: Oxford/Mikolajczyk,
HPatches, plus our blob-marked-skin captures and synthetic warps with known
strain. Metrics: latency/frame, throughput (Mpix/s), DRAM bytes (Nsight),
occupancy; repeatability/matching score; strain RMSE. (See Option~3 harness.)
\end{fillbox}

%==============================================================================
\section{Results}
\label{sec:results}

\subsection{Throughput and Latency}
\begin{fillbox}
Table + plot: ms/frame and Mpix/s across resolutions (VGA$\rightarrow$4K) vs.\
baselines. Headline speedup.
\end{fillbox}

\subsection{Memory Traffic and the Fusion Win}
\begin{fillbox}
Nsight DRAM bytes: fused DoG+extrema vs.\ materialized. Roofline placement.
\end{fillbox}

\subsection{Detection Quality and Repeatability}
\begin{fillbox}
Repeatability under rotation/scale/illumination; keypoint-count parity so speed
isn't bought with worse features.
\end{fillbox}

\subsection{Strain-Field Accuracy}
\begin{fillbox}
Recovered strain vs.\ ground truth on flexing-joint captures; error maps.
\end{fillbox}

\subsection{Ablations}
\begin{fillbox}
(a) fused vs.\ materialized DoG; (b) fused col\_row vs.\ two-pass;
(c) persistent buffers vs.\ per-frame alloc; (d) CoV threshold sweep;
(e) \texttt{maxregs}/block-size design space.
\end{fillbox}

%==============================================================================
\section{Discussion and Limitations}
\label{sec:discussion}
\begin{fillbox}
Radial-flux signature gives symmetry-based selection but no canonical rotation
like Lowe's---state clearly. Shared-memory capacity bounds patch radius / layer
count. Behavior under extreme scale change and marker occlusion.
\end{fillbox}

%==============================================================================
\section{Conclusion}
\label{sec:conclusion}
\begin{fillbox}
Recap the fully-resident design, the novel signature/filter, and the strain
application. Future work: multi-GPU, descriptor extension, on-device matching.
\end{fillbox}

%------------------------------------------------------------------------------
\begin{thebibliography}{1}
\bibitem{lowe2004} D.~G. Lowe, ``Distinctive image features from scale-invariant keypoints,'' \emph{IJCV}, 2004.
\bibitem{surf} H.~Bay et al., ``SURF: Speeded up robust features,'' \emph{CVIU}, 2008.
\bibitem{kaze} P.~Alcantarilla et al., ``KAZE features,'' \emph{ECCV}, 2012.
\bibitem{mikolajczyk2005} K.~Mikolajczyk et al., ``A comparison of affine region detectors,'' \emph{IJCV}, 2005.
\bibitem{siftgpu} C.~Wu, ``SiftGPU: A GPU implementation of SIFT,'' 2007.
\bibitem{popsift} C.~Griwodz et al., ``PopSift: A faithful SIFT implementation for real-time applications,'' \emph{MMSys}, 2018.
\bibitem{cudasift} M.~Bjorkman, ``CudaSift,'' \note{cite/URL}.
\bibitem{halide} J.~Ragan-Kelley et al., ``Halide,'' \emph{PLDI}, 2013.
\bibitem{micikevicius_stencil} P.~Micikevicius, ``3D finite difference computation on GPUs using CUDA,'' \emph{GPGPU}, 2009.
\bibitem{cuda_warpagg} \note{NVIDIA dev blog, warp-aggregated atomics / ballot compaction.}
\bibitem{dic} M.~A. Sutton et al., \emph{Image Correlation for Shape, Motion and Deformation Measurements}, Springer, 2009.
\end{thebibliography}
